\documentclass{article}
\usepackage{amsmath}
\usepackage{amssymb}
\usepackage{amsthm}
\usepackage{MnSymbol}
\usepackage{xcolor}
\usepackage{parskip}
\usepackage{tikz}
\usepackage{bbm}
\usepackage{hyperref}
\usetikzlibrary{trees}
\usetikzlibrary{graphs}


\theoremstyle{definition}
\newtheorem{definition}{Definition}
\newtheorem{theorem}{Theorem}


\newcommand{\abs}[1]{\left| #1 \right|}


\title{lecture}
\date{\today}
\begin{document}
\maketitle


ranking, correlation clustering, hierarchical clustering

for 2 items, ask oracle for ordering. similarly for 3, ask for which one is outside, for 4, ask for which 2 are grouped.

2 has 1 wrong answer, 3 and 4 have 2 wrong answers


convert hte 2 item situation into a directed graph like from l7. \(a<b\) implies directed edge from a to b that is positive and a negative from b to a

if \(m\) represents the number of constraints, \(m_s\) are the number of constraints satisfied, \(m_v\) the number violated, and \(m_a\) unaffected, we see that


\begin{equation}
    W(L,r) = m_s-m_v
\end{equation}

\begin{equation}
    ALG = m_s + \frac{1}{2} \left( m-m_s-m_v \right) = \frac{1}{2}(m+m_s-m_v) = \frac{m}{2}+\frac{1}{2}W(L,R)
\end{equation}



since the weights are bounded, we can increase the weights uniformly (?)

but after increasing, we have to undo the inrease which results in a minus total weight term. 


for triplets, create a graph that is like a triangle, if bc are together, the edge is -2, ab and ac are +1. 

\begin{equation}
    W(L,r) = 2m_s-m_v
\end{equation}

\begin{equation}
    ALG = m_s + \frac{1}{3}(m-m_s-m_v) = \frac{1}{3}(m+w(L,r))
\end{equation}

for quartets, if ab and cd are together, connect them with -2, then all other configs with +1



\end{document}